\chapter{Cái Bẫy Của Món Nhỏ}
\markboth{Cái Bẫy Của Món Nhỏ}{The Trap of Small Purchases}

\section{Tháng Không Mua Gì Lớn Vẫn Hết Tiền}
\begin{truyen}{Những Khoản Không Đáng Bao Nhiêu}{The Expenses That Seem Insignificant}
\chuhoa{C}{ó một tháng người thầy không mua bất cứ món đồ lớn nào. Không tivi mới, không điện thoại mới, không chuyến đi đắt tiền nào. Thế mà đến cuối tháng ví vẫn rỗng.}
\textit{(There was a month when the teacher bought nothing major. No new television, no new phone, no expensive trip. Yet by the end of the month his wallet was still empty.)}

Ngồi nhớ lại, anh mới thấy tiền đi qua những chỗ rất nhỏ: một ly cà phê cho tỉnh, một bữa ăn ngoài vì ngại nấu, vài món giảm giá vì tưởng rằng rẻ thì mua cũng không sao. Mỗi khoản một chút, nhẹ như bụi. Nhưng bụi bám đủ nhiều thì cũng phủ mờ cả mặt bàn.
\textit{(As he looked back, he saw money leaving in very small ways: a coffee for alertness, a meal outside because cooking felt tiring, a discounted item because cheap seemed harmless. Each expense was tiny, almost like dust. Yet when enough dust gathers, it can cover an entire table.)}

Điều nguy nhất của món nhỏ là nó đi qua mà không để lại cảm giác đau. Người ta chỉ đau khi rút một khoản lớn. Còn những thứ nhỏ lặp đi lặp lại lại mòn ví một cách lặng lẽ. Khi nhận ra, người ta thường đã quen tay mất rồi.
\textit{(The most dangerous thing about small purchases is that they pass without pain. People usually feel pain only when paying a large amount. But repeated little things wear down the wallet quietly. By the time one notices, the hand is often already accustomed to them.)}
\end{truyen}

\begin{baihoc}
Lãng phí hiếm khi đến trong hình dạng ồn ào. Nó thường đến bằng những điều rất nhỏ nhưng lặp lại rất đều.
\textit{(Waste rarely arrives loudly. It usually comes in very small things repeated with great consistency.)}
\end{baihoc}

\begin{ghinhoanh}
Small leaks sink strong boats.

Quiet habits can become costly.
\end{ghinhoanh}

\begin{tuhocnhanh}
1. Ba khoản nhỏ nào đang lặp lại nhiều nhất trong tháng của em? \textit{Which three small expenses repeat most in your month?}

2. Em đã bao giờ mất tiền vì nghĩ rằng ``chỉ có một chút'' chưa? \textit{Have you ever lost money because you thought, ``it is only a little''?}
\end{tuhocnhanh}
\ngancach

\section{Một Hạt Gạo Một Sợi Chỉ}
\begin{truyen}{Lời Dạy Về Nguồn Gốc Của Vật}{A Teaching About the Origin of Things}
\chuhoa{N}{gười xưa có câu: ``Nhất châu nhất phạn, đương tư lai xứ bất dị; bán ti bán lũ, hằng niệm vật lực duy gian.'' Một hạt gạo, một sợi chỉ đều không dễ mà có.}
\textit{(The ancients said, ``For each grain of rice and every thread of cloth, remember how difficult their origins are.'' A single grain of rice, a single thread, never comes easily.)}

Điều hay của lời dạy này là nó đưa con người trở về với gốc của vật dùng hằng ngày. Khi thấy một món đồ, thay vì chỉ nhìn giá tiền, ta bắt đầu nhớ đến thời gian, công sức, mồ hôi của bao người phía sau. Cách nhìn ấy làm lòng người bớt tùy tiện.
\textit{(The beauty of this teaching is that it brings a person back to the origin of everyday things. When seeing an item, instead of looking only at its price, we begin to remember the time, labor, and sweat of many people behind it. Such a perspective makes the heart less careless.)}

Biết ơn nguồn gốc của vật là một cách chống lãng phí rất sâu. Bởi ai nhớ bát cơm khó nhọc mà có sẽ khó nỡ bỏ thừa. Ai nhớ đồng tiền vất vả mà có sẽ khó rút ra vô tâm.
\textit{(Gratitude for the origin of things is a deep way to resist waste. Whoever remembers how hard a bowl of rice was earned will struggle to throw it away. Whoever remembers how hard money was earned will hesitate to spend it carelessly.)}
\end{truyen}

\begin{baihoc}
Muốn quý tiền, hãy học cách quý những gì tiền đại diện: thời gian, công sức và sự nhọc nhằn.
\textit{(To value money, learn to value what money represents: time, effort, and hardship.)}
\end{baihoc}

\begin{ghinhoanh}
Value the labor behind the object.

Gratitude reduces waste.
\end{ghinhoanh}

\begin{tuhocnhanh}
1. Em đang dùng thứ gì mỗi ngày mà đã quên mất cái giá của nó? \textit{What do you use every day whose true cost you have forgotten?}

2. Biết ơn có làm em tiêu tiền chậm lại không? \textit{Does gratitude slow down your spending?}
\end{tuhocnhanh}
\ngancach

\section{Ghi Chép Để Nhìn Ra Vết Rò}
\begin{truyen}{Cuốn Sổ Mỏng Nhưng Có Tác Dụng Lớn}{A Thin Notebook with Great Power}
\chuhoa{S}{au tháng ấy, người thầy mua một cuốn sổ nhỏ. Mỗi tối anh ghi lại tất cả những gì đã chi trong ngày, kể cả những khoản rất nhỏ.}
\textit{(After that month, the teacher bought a small notebook. Every evening he recorded everything he had spent during the day, including very small amounts.)}

Ban đầu việc ghi chép khiến anh thấy phiền. Nhưng chỉ sau vài tuần, cuốn sổ trở thành tấm gương. Nó không mắng mỏ, không làm nhục ai, chỉ phản chiếu rất rõ rằng tiền đang chảy đi đâu. Có những thứ khi còn nằm trong đầu thì mơ hồ, nhưng khi đã viết ra trên giấy thì hết đường tự bào chữa.
\textit{(At first the record-keeping felt troublesome. But after a few weeks, the notebook became a mirror. It did not scold or shame anyone; it simply reflected very clearly where the money was going. Some things remain vague while in the mind, but once written on paper there is no room left for excuses.)}

Từ đó, anh hiểu rằng quản tiền bắt đầu từ việc dám nhìn thẳng. Người không nhìn số tiền mình tiêu mỗi ngày thường cũng là người ngạc nhiên nhất khi thấy mình không còn gì cuối tháng.
\textit{(From then on, he understood that managing money begins with the courage to look directly at it. Those who never look at what they spend each day are often the most surprised when nothing is left at month’s end.)}
\end{truyen}

\begin{baihoc}
Ghi chép không tạo ra tiền, nhưng nó ngăn tiền biến mất trong sương mù của thói quen.
\textit{(Record-keeping does not create money, but it prevents money from disappearing into the fog of habit.)}
\end{baihoc}

\begin{ghinhoanh}
Write it down.

What is visible is easier to correct.
\end{ghinhoanh}

\begin{tuhocnhanh}
1. Em có dám ghi lại toàn bộ chi tiêu trong bảy ngày tới không? \textit{Would you dare record all spending for the next seven days?}

2. Điều gì em sợ nhất khi phải nhìn con số thật? \textit{What do you fear most about seeing the real numbers?}
\end{tuhocnhanh}

\begin{turanminh}
1. Tôi phải nhớ rằng chính những khoản ``không đáng bao nhiêu'' đã nhiều lần làm cuối tháng của mình nặng nề.

2. Tôi không được coi thường cà phê lặt vặt, món rẻ mua cho vui hay những khoản tiện tay, vì nhiều lỗ rò nhỏ cộng lại có thể đánh bại cả một tháng lao động.

3. Mỗi tối ghi chép vài dòng tuy nhỏ, nhưng đó là cách tôi tự kéo mình ra khỏi sự mơ hồ của thói quen cũ.
\end{turanminh}

\begin{dayhoc}
1. Có thể giao học sinh một bài tập ngắn: ghi ba khoản nhỏ các em hay xin cha mẹ và tự phân tích khoản nào là cần, khoản nào chỉ là thích.

2. Nhắc học sinh rằng lãng phí thường không bắt đầu từ việc lớn, mà từ thói quen xem nhẹ những điều nhỏ.

3. Một câu gợi mở tốt trên lớp là: ``Em có đang làm cha mẹ mất tiền vì những thứ em nghĩ là rất nhỏ không?''
\end{dayhoc}